\section{Lecture — 01/10/26}


\subsection*{Curriculum Position}
\begin{itemize}
    \item Week-2 — Forces, Vectors, and Vector Operations
\end{itemize}

\subsection*{Key Content}
\begin{enumerate}
    \item Vector Addition of Forces
    \item Cartesian Vector Notation and the Decomposition of Vectors
    \item Right Handed Coordinate Systems and 3D Vector Decomposition
    \item Unit Vectors and the Direction of a Cartesian Vector
    \item Concurrent Force Systems and x, y, z Coordinates 
    \item Position Vectors and Force Vectors Directed Along a Line
    \item Dot Products, Cartesian Vector Formulation
    
    
\end{enumerate}

\subsection*{Instructor Emphasis}
\begin{itemize}
    \item Working with units while solving problems is highly encouraged \\ 
        for unit consistency and orders of magnitude assurance.
    \item The instructor consistently preferred the laws of sines, cosines, and parallelograms      \\ over vector decomposition when it came to simpler problems. 
    \item Vector decomposition is advised for the addition of more than 2 vectors.
    \item Studying example 2.17 is recommended.   
        
\end{itemize} 

\subsection*{Announcements}


\subsection*{Open Questions}


\subsection*{Independent Follow-up}
\begin{itemize}
    \item Minimum: Work on the vector addition examples provided in the lecture notes.
    \item Exploration: Watch some stuff on statics and FEA.
\end{itemize}
